\section*{第 \ref{ch:g7:moon-phases} 章——月相}

\begin{solution}{exo:g7:moon-phases:1}
永远恰好一半——朝向太阳的那一半。\omterm{def:g7:moon-phases:phases}{月相}只是
观察的角度：被照亮的那一半中恰好朝向
地球的那部分所占的比例。
\end{solution}

\begin{solution}{exo:g7:moon-phases:2}
依次为新月、\omterm{def:g7:moon-phases:phases}{渐盈}\omterm{ex:g2:moon-in-the-sky:shapes}{蛾眉月}、上弦月、\omterm{def:g7:moon-phases:phases}{渐盈}凸月、\omterm{ex:g2:moon-in-the-sky:shapes}{满月}、
\omterm{def:g7:moon-phases:phases}{渐亏}凸月、下弦月、\omterm{def:g7:moon-phases:phases}{渐亏}\omterm{ex:g2:moon-in-the-sky:shapes}{蛾眉月}——循环一周约 $29.5$ 天。
\omterm{def:g7:moon-phases:phases}{渐盈}是指从新月到\omterm{ex:g2:moon-in-the-sky:shapes}{满月}的增亮阶段。
\end{solution}

\begin{solution}{exo:g7:moon-phases:3}
新月：你面对灯，小球朝向灯——
它被照亮的一面背对你。\omterm{ex:g2:moon-in-the-sky:shapes}{满月}：灯在你身后，小球在你面前——
它被照亮的一面完全朝向你。
\end{solution}

\begin{solution}{exo:g7:moon-phases:4}
任选两条：地球的\omterm{def:g1:light-and-shadows:shadow}{影子}背向太阳，只可能落到
处于\omterm{ex:g2:moon-in-the-sky:shapes}{满月}位置的\omterm{def:g2:moon-in-the-sky:moon}{月球}上；\omterm{ex:g2:moon-in-the-sky:shapes}{蛾眉月}位于我们朝向太阳的一侧，远离
那个\omterm{def:g1:light-and-shadows:shadow}{影子}；圆形\omterm{def:g1:light-and-shadows:shadow}{影子}的边缘永远无法
勾勒出凸月那样的双重曲线。
\end{solution}

\begin{solution}{exo:g7:moon-phases:5}
新月对我们来说是暗的，因为它的\emph{自身}夜面
朝向我们——与我们的\omterm{def:g1:light-and-shadows:shadow}{影子}无关，远侧则处于
阳光的完全照射之下。\omterm{def:g7:shadows-eclipses:lunar}{月食}是\omterm{ex:g2:moon-in-the-sky:shapes}{满月}穿过地球的
真正\omterm{prop:g7:shadows-eclipses:umbra}{本影}：我们的\omterm{def:g1:light-and-shadows:shadow}{影子}确实短暂地落在\omterm{def:g2:moon-in-the-sky:moon}{月球}上。
\end{solution}

\begin{solution}{exo:g7:moon-phases:6}
\omterm{ex:g2:moon-in-the-sky:shapes}{满月}位于太阳的对面。随着旋转的地球带着太阳
沉入西方地平线以下，天空中与之相对的那一点——
即\omterm{ex:g2:moon-in-the-sky:shapes}{满月}所在之处——必定正在东方升起：
\omterm{def:g1:day-and-night:daynight}{日落}与\omterm{ex:g2:moon-in-the-sky:shapes}{满月}升起，
是同一事件的两种看法。
\end{solution}

\begin{solution}{exo:g7:moon-phases:7}
新月与太阳同行：午夜时，二者都在你脚下，
位于地球的另一侧。午夜时见到新月，
与午夜时见到太阳一样不可能。
\end{solution}

\begin{solution}{exo:g7:moon-phases:8}
远离\omterm{def:g1:day-and-night:daynight}{日落}点——\omterm{ex:g2:moon-in-the-sky:shapes}{蛾眉月}的弧线朝向太阳，
两角背离太阳。被照亮的一面必须朝向太阳，
两角是被照亮一面的后退边缘，
几何关系使它们不可能指向太阳一侧。
\end{solution}

\begin{solution}{exo:g7:moon-phases:9}
介于上弦月与\omterm{ex:g2:moon-in-the-sky:shapes}{满月}之间——大约绕行了全程的八分之三，
位于傍晚一侧。它在午后升起。五晚之后：\omterm{ex:g2:moon-in-the-sky:shapes}{满月}——
\omterm{def:g1:day-and-night:daynight}{日落}时升起，整夜都在空中。
\end{solution}

\begin{solution}{exo:g7:moon-phases:10}
\omterm{def:g7:shadows-eclipses:solar}{日食}：只发生在新月时——太阳–\omterm{def:g2:moon-in-the-sky:moon}{月球}–地球排成一线。
\omterm{def:g7:shadows-eclipses:lunar}{月食}：只发生在\omterm{ex:g2:moon-in-the-sky:shapes}{满月}时——太阳–地球–\omterm{def:g2:moon-in-the-sky:moon}{月球}排成一线。即便如此也很少见，
因为\omterm{def:g2:moon-in-the-sky:moon}{月球}倾斜的\omterm{def:g6:sun-earth-moon:axisorbit}{轨道}通常使它位于
准确直线的上方或下方：\omterm{def:g7:moon-phases:phases}{月相}是必要条件，
排成一线却无法保证。
\end{solution}

\begin{solution}{exo:g7:moon-phases:11}
地球的\omterm{def:g7:moon-phases:phases}{月相}与我们所见相反，且处在同一时刻：
当我们看到新月时，\omterm{def:g2:moon-in-the-sky:moon}{月球}宇航员看到的是\emph{满地球}——
我们整个被照亮的一面朝向他们；我们看到\omterm{ex:g2:moon-in-the-sky:shapes}{满月}时，他们看到的是新地球。
同一种几何关系，换一个座位来看。
\end{solution}

\begin{solution}{exo:g7:moon-phases:12}
路径：阳光照射地球，明亮的地球将其向\omterm{def:g2:moon-in-the-sky:moon}{月球}漫射，
\omterm{def:g2:moon-in-the-sky:moon}{月球}的夜面再把少量光向回漫射到我们眼中——阳光被转送了两次。
在我们看到\omterm{ex:g2:moon-in-the-sky:shapes}{蛾眉月}相时，\omterm{def:g2:moon-in-the-sky:moon}{月球}看到的几乎是
\emph{满地球}（即相反\omterm{def:g7:moon-phases:phases}{月相}规律）：
\omterm{def:g2:moon-in-the-sky:moon}{月球}的夜面被尽可能明亮的地球照亮——
地照最强。临近我们的\omterm{ex:g2:moon-in-the-sky:shapes}{满月}时，\omterm{def:g2:moon-in-the-sky:moon}{月球}看到的是细窄的
蛾眉形地球：照明变弱，满\omterm{def:g7:moon-phases:phases}{月相}的强光
淹没了残余的微弱地照。
\end{solution}

\begin{solution}{pb:g7:moon-phases:1}
\textbf{1.} 上弦月约在第 $7$ 夜；\omterm{ex:g2:moon-in-the-sky:shapes}{满月}约在第
$15$ 夜；下弦月约在第 $22$ 夜（周期约在第
$29$--$30$ 夜前后结束）。
\textbf{2.} \omterm{ex:g2:moon-in-the-sky:shapes}{满月}：\omterm{def:g1:day-and-night:daynight}{日落}时升起。上弦月：
约在正午升起。下弦月：约在午夜升起。新月：
与太阳同升同落。
\textbf{3.} 初月在\omterm{def:g6:sun-earth-moon:axisorbit}{轨道}上仅稍稍落后于太阳，
\omterm{def:g1:day-and-night:daynight}{日落}后挂在西方，很快随太阳落下。\omterm{ex:g2:moon-in-the-sky:shapes}{满月}位于对面：
太阳在西方落下时，
对面则意味着在东方升起。
\textbf{4.} 任何与太阳的移动部分相同的\omterm{def:g7:moon-phases:phases}{月相}：\omterm{def:g7:moon-phases:phases}{渐盈}凸月
整个下午可见，尤其是\omterm{def:g7:moon-phases:phases}{渐亏}的各阶段——下弦月
与\omterm{def:g7:moon-phases:phases}{渐亏}\omterm{ex:g2:moon-in-the-sky:shapes}{蛾眉月}
在清晨的天空中高挂。
\textbf{5.} \omterm{def:g7:moon-phases:phases}{渐盈}凸月：小球被照亮四分之三，说明你位于
上弦月与\omterm{ex:g2:moon-in-the-sky:shapes}{满月}之间——灯在你的肩后，
小球在你的前方偏侧。
\textbf{6.} 他们看到的是小球\emph{自身}的夜面——
背对灯光的那一半。你所在的地球也会投下\omterm{def:g1:light-and-shadows:shadow}{影子}，
但它向你身后延伸，远离灯光；小球只有
在\omterm{ex:g2:moon-in-the-sky:shapes}{满月}位置时才会落入其中，
而且还须恰好排成一线：那是\omterm{def:g7:shadows-eclipses:lunar}{月食}，而不是\omterm{def:g7:moon-phases:phases}{月相}。
\textbf{7.} \omterm{def:g2:moon-in-the-sky:moon}{月球}恰好每绕行一周便自转一周，因此同一
半球永远朝向地球——连玉兔的图案在内；背面
则直到有摄像机绕飞过去才被看到。
\textbf{8.} 自转与公转保持完美同步：\omterm{def:g2:moon-in-the-sky:moon}{月球}沿\omterm{def:g6:sun-earth-moon:axisorbit}{轨道}前进时，
它恰好转过足够的角度，使同一面始终朝向我们——
就像舞者绕着舞伴转圈时始终面朝对方：
确实在自转（旁观者能看到各个侧面），
却只向舞伴展示一面。
\textbf{9.} 午夜升起对应下弦月的时刻表——
即\omterm{def:g7:moon-phases:phases}{渐亏}的一侧：下弦月（被照亮的一半
朝向清晨太阳的方向）。
\textbf{10.} 错误之处：\omterm{ex:g2:moon-in-the-sky:shapes}{蛾眉月}靠近太阳，因此午夜时
绝不可能高挂空中；两角\emph{背离}太阳，从不指向太阳；
午夜出现\omterm{def:g1:day-and-night:daynight}{日落}色彩，本身同样不可能；
三次出错，三振出局。
\textbf{11.} 不能：\omterm{def:g7:shadows-eclipses:solar}{日食}要求\omterm{def:g2:moon-in-the-sky:moon}{月球}位于太阳与地球之间——
即新月的位置。今晚\omterm{def:g2:moon-in-the-sky:moon}{月球}处于\omterm{ex:g2:moon-in-the-sky:shapes}{满月}，位于相对的一侧：
两件事不可能
同在一夜发生。
\textbf{12.} 例如：“一个事实贯穿整个月：\omterm{def:g2:moon-in-the-sky:moon}{月球}总有一半被照亮，
我们从不断变换的位置观察这一半。本月还要留意：地球的\omterm{def:g1:light-and-shadows:shadow}{影子}
与\omterm{def:g7:moon-phases:phases}{月相}毫无关系，它只在罕见的\omterm{def:g7:shadows-eclipses:lunar}{月食}中出现。
还要记住这份时刻表：每种\omterm{def:g7:moon-phases:phases}{月相}都有自己的升起时刻，
\omterm{ex:g2:moon-in-the-sky:shapes}{满月}在\omterm{def:g1:day-and-night:daynight}{日落}时升起，两个弦月分别在正午和午夜升起，
因此\omterm{def:g2:moon-in-the-sky:moon}{月球}的形状
便能告诉你它的时刻安排。”
\end{solution}
